\documentclass{article}
\usepackage[T1]{fontenc}
\usepackage{lmodern}
\usepackage{graphicx}
\usepackage{amsmath,amssymb}
\usepackage[a4paper,margin=2cm]{geometry}
\usepackage[absolute,overlay]{textpos}
\setlength{\TPHorizModule}{1mm}
\setlength{\TPVertModule}{1mm}
\usepackage[indonesian]{babel}
\usepackage{hyperref}
\setlength{\emergencystretch}{2em}
% unit: Halaman 14

\begin{document}

% === Halaman 14 (python/native) ===

14

• OTP ini tidak dapat dipecahkan  karena:

   1. Kunci OTP (acak) + plainteks (tidak  acak) = cipherteks (acak)


2. Hanya terdapat satu kunci yang memetakan plainteks ke

        cipherteks, begitu juga sebaliknya 

       → fungsi C = (P + K) mod 26 adalah pemetaan satu-ke-satu (one-to-one).

       → artinya, untuk sembarang plainteks dan cipherteks hanya ada satu 

kunci yang memetakannya satu sama lain

       → Akibatnya, jika kriptanalis mencoba mendekripsi cipherteks dengan 

beberapa  kunci berbeda ada kemungkinan menghasilkan beberapa 

plainteks yang bermakna, sehingga kriptanalis kesulitan menentukan 

plainteks mana yang benar.

Enkripsi: C = (P + K) mod 26

Dekripsi: P = (C – K) mod 26

\end{document}
